\section{Three Obligations in One Sentence}

Canon 277 §1 of the Latin Code imposes two things and calls them one. \latin{Clerici obligatione tenentur servandi perfectam perpetuamque propter Regnum coelorum continentiam, ideoque ad coelibatum adstringuntur}---clerics are held by the obligation of observing perfect and perpetual continence for the sake of the Kingdom of heaven, and therefore are bound to celibacy.\endnote{CIC c.~277 §1, quoted from the Holy See's Latin web delivery of Book~II, fetched, hashed, and read 2026-07-25 (see References; the exact response was byte-identical to the state already registered in this repository's source library). The delivery prints \latin{coelorum} and \latin{coelibatum} with the \emph{oe} digraph where the \emph{Acta Apostolicae Sedis} promulgation text of the Eastern Code prints \latin{caelibatus}; orthographic variation between and within official Latin texts is normal and carries no juridical weight. Where this article quotes an English rendering of a Latin canon inside quotation marks, the rendering is the Holy See's own English web delivery, identified in the References; where no quotation marks appear, the English is the article's own working gloss and the Latin governs.} The canon has two objects and one \latin{ideoque}. Continence---abstention from sexual activity---is what the cleric is obliged to observe. Celibacy---the unmarried state---is what he is therefore bound to. The English delivery of the same canon says it plainly: clerics ``are obliged to observe perfect and perpetual continence for the sake of the kingdom of heaven and therefore are bound to celibacy.''

Twenty-two canons later, in the law for consecrated life, the Code sets the same two words under a third. Canon 599: \latin{Evangelicum castitatis consilium propter Regnum coelorum assumptum \ldots\ obligationem secumfert continentiae perfectae in caelibatu}---the evangelical counsel of chastity, assumed for the sake of the Kingdom of heaven, ``entails the obligation of perfect continence in celibacy.''\endnote{CIC c.~599, Latin delivery of Book~II; English from the Vatican delivery of cann.~573--606, fetched, hashed, and read 2026-07-25. The canon is the Code's most compact statement of the relation among the three terms: chastity is the counsel; continence is what the counsel obliges; celibacy is the state in which that continence is lived.} Chastity is the counsel; continence is its obligation; celibacy is the condition in which the obligation is lived. Three words, arranged in a hierarchy, none of them a synonym for another.

And in the Eastern Code, promulgated seven years later for the twenty-three Eastern Catholic Churches, canon 374 puts one of the three words---and only one---on every cleric alike: \latin{Clerici caelibes et coniugati castitatis decore elucere debent}. Clerics, celibate \emph{and married}, must shine with the beauty of chastity.\endnote{CCEO c.~374, quoted from the promulgation text printed in \emph{Acta Apostolicae Sedis} 82 (1990) 1142, read at the page image of that page in the Holy See's archival PDF of the volume, fetched and hashed 2026-07-25 (see References). Every CCEO canon quoted in this article was read at its \emph{Acta} page image, not merely at the volume's OCR layer; the OCR served as a finding aid. English renderings of CCEO canons throughout are the article's own working glosses: the Code has no official English version, and the Latin alone is authentic.} The same sentence that binds a celibate Eastern presbyter binds a married one. It says nothing about continence and nothing about celibacy, because in that Code those are not what the word means.

Ordinary speech collapses all three into one. ``Celibacy'' is made to carry the moral virtue, the physical abstention, and the juridical status at once, so that a married Ukrainian Greek Catholic priest is described as an exception to celibacy---which he is---and, in the same breath, as an exception to chastity, which he is not and which no Catholic law has ever suggested. The confusion is not merely verbal. It reliably damages both sides of the argument that runs through this subject: a defence of the discipline that argues from the value of chastity has defended something no one disputes, and a critique that treats the obligation as a single indivisible package has attacked something the law does not contain. The law itself is more finely jointed than the argument about it.

This article works those joints. It asks what the Catholic Church's current law obliges of her clerics in this matter; on what authority; with what effects when the obligation is broken; how the two codes differ and why the difference is not an exception to a rule; and what the history and the magisterium do and do not establish. The order is the law's own. Sections~2 through~6 fix the three terms and take the Latin discipline apart canon by canon---the obligation, the acts by which it attaches to a man, the marriage impediment and penal norm behind it, and the two distinct acts by which it is undone. Section~7 turns to the Eastern common law at its promulgation text; section~8 sets the twenty-four Churches \latin{sui iuris} side by side; section~9 takes up the permanent diaconate, where the Latin Code's vocabulary produces its sharpest unsettled question. Section~10 reconstructs how the Western discipline was built, at identified witnesses; sections~11 and~12 cover convert clergy and the modern magisterial file; section~13 states the theological arguments on both sides at their strongest without adjudicating what the Church has not.

Two boundaries govern everything that follows. First, almost all of this is mutable law. It is verified against the competent current texts as of 25 July 2026, and the Legal Scope and Currentness statement in the terminal appendix records exactly which bodies of law, which editions, which jurisdictions, and which amendments that verification covers. Second, no concrete case is decided here. Whether a particular man is bound, dispensed, impeded, penalized, or free to marry is a question of fact and of competent authority, and section~14 says where such a question belongs. This is a study of norms, written for readers who want to know what the norms say.
